\section{Familias de Optimizadores y Dinámicas de Entrenamiento}
\label{sec:optimizer-families}

La optimización de arquitecturas transformer altamente parametrizadas presenta desafíos significativos debido a paisajes de pérdida no convexos, puntos de silla y heterogeneidad de bloques entre grupos de parámetros. Este sistema aborda estos desafíos mediante un marco unificado que soporta veintitrés familias de optimizadores que abarcan seis categorías algorítmicas: (1) líneas base clásicas (SGD, AdamW), (2) momento avanzado y reducción de varianza (Adan, ADOPT, AdEMAMix, MARS, Cautious), (3) variantes eficientes en memoria (Adafactor, GaLore, Lion, APOLLO, APOLLO-Mini, Q-APOLLO), (4) métodos sin programación y sin parámetros (Schedule-Free AdamW, Prodigy) más escalado de lotes grandes mediante LAMB, (5) conscientes de curvatura y segundo orden (Sophia, Shampoo, SOAP), (6) orientados a geometría (Muon, Turbo-Muon), y (7) optimizadores con adaptabilidad ajustable (Anon). Las descripciones algorítmicas detalladas, pseudocódigo y una tabla comparativa de memoria/complejidad para todos los optimizadores se proporcionan en el Apéndice~\ref{sec:annex-optimizers}.

\begin{figure}[t]
\centering
\fitfigure{\begin{tikzpicture}[node distance=1.1cm and 1.0cm]
\node[box, fill=blue!8] (start) {Elegir objetivo de optimizador};
\node[box, fill=green!8, below left=of start, xshift=-1.5cm] (baseline) {Línea base confiable\\AdamW / RAdam};
\node[box, fill=yellow!12, below=of start] (memory) {Menor memoria de optimizador\\Adafactor / GaLore / Lion};
\node[box, fill=orange!12, below right=of start, xshift=1.5cm] (advanced) {Agresivo o estructurado\\Adan, ADOPT, Sophia,\\Shampoo, SOAP, Muon};
\node[box, fill=red!8, below=of memory] (route) {Grupos con prefijo de esquema\\LR, decay de peso, betas, eps};
\draw[line] (start) -- (baseline);
\draw[line] (start) -- (memory);
\draw[line] (start) -- (advanced);
\draw[line] (baseline) -- (route);
\draw[line] (memory) -- (route);
\draw[line] (advanced) -- (route);
\end{tikzpicture}}
\caption{Selección de optimizador en la práctica: la clase determina la regla de actualización, mientras que el esquema controla cómo se aplican los hiperparámetros entre embeddings, normas, bloques recurrentes, bloques de atención y otros parámetros.}
\label{fig:optimizer-routing}
\end{figure}

